\documentclass[UTF8]{article}
\usepackage{amsmath}
\usepackage{amssymb}
\usepackage{bm}
\usepackage{geometry}
\geometry{a4paper, margin=1in}

\title{\textbf{ZGFT (Ziqi Geometry Field Theory) Introduction: On the Four-Dimensional Form, N+a May Be a}}
\author{Ziqi Zhi}
\date{v4.8}

\begin{document}
\maketitle

\section{Definitions}

Oh, by the way, because of possible rendering errors of the language, all \(\mathbf{k}\) below are in bold; if \(\mathbf{k}\) carries a suffix or a power subscript, then it is in white body. Bold \(\mathbf{k}\) is a vector.

\begin{equation}
\Delta_+ f(x)=\frac{f(x+a)-f(x)}{a}
\end{equation}
(This is microscopic; if it is macroscopic, then this and all the following are central differences.)

At the same time, this formula and all the following default to \(c=1\).

\begin{equation}
\Delta_{++}(x)=\frac{f(x+2a)-2f(x+a)+f(x)}{a^2}
\end{equation}

\begin{equation}
\Delta_{2x+}f(x,y)=\frac{f(x+a,y)-f(x,y)}{a}
\end{equation}

\(\Delta_{2y+}\) is the same.

\(\Delta_{2x++}\) and all the rest are also omitted.

\(\sum f(x)a\) is a discrete integral; the first and second laws of calculus still hold.

When \(a\to 0\), it degenerates into traditional calculus.

Physically it is impossible for it to take 0; the minimum distance can only be the Planck distance, after all this is the smallest meaningful place.

On the four-dimensional form \(N+a\) may be \(a\). Therefore, there is no boundary. For example, on a sphere, we walk a little more than one circle and we only return to the starting point; and spacetime as four dimensions will be more complicated, therefore, there is no boundary.

\section{ZGFT Part}

\subsection{Discrete variables}

Define \(\mathbf{k}\) (the discrete variable):

\(k_x\) is the change of \(a\) on the \(x\) axis, \(k_y\) is the change of \(a\) on the \(y\) axis. \(k_j\in\mathbb{C}\) (complex numbers), where \(j=(x,y,\dots)\)

\(\mathbf{k}^4\) is the combined change of \(\mathbf{k}\) in the four directions. It can directly replace the four directions.

\(\mathbf{k}^4\) can be written in the form of a matrix:

Assume

\begin{align*}
k_x &= 6+5k_y+2k_z\\
k_y &= 3+5k_t\\
k_z &= 2\\
k_t &= 1
\end{align*}

\begin{equation}
\mathbf{k}^4=
\begin{pmatrix}
6 & 5 & 2 & 0\\
0 & 3 & 0 & 5\\
0 & 0 & 2 & 0\\
0 & 0 & 0 & 1
\end{pmatrix}
\end{equation}

Let it be stated once more, in the first column, the first row, that \(6\) does not mean \(k_x\) times \(6\), but \(k_x=6\). Every main diagonal is like this. The value on the diagonal represents its \(k\) value, and the other elements not on the main diagonal explain their coupling by the way of the column. Other rows are the same.

And the normal way to find the small volume of a region of an object is \(a^3\); if time is also counted, then multiply also by \(a_tc\).

\(a_t\) should generally be the Planck time.

Its two-dimensional Pythagorean law expression is:

\begin{equation}
S^2=\operatorname{Re}\left[\left(\sum_{n}^{x_1}k_{xn}a\right)^2+\left(\sum_{n}^{y_1}k_{yn}a\right)^2\right]
\end{equation}

The left side and the right side of this formula are both real numbers.

Here \(x_1\) is the total distance on \(x\), \(y_1\) is the same.

One can also apply \(\Delta_+\) to both sides of this formula.

The minimum distance between two points is \(a\); the minimum volume and area are \(a^3\) and \(a^2\).

\subsection{Gravity}

\begin{equation}
\Delta_{++}(Z)=\sum_{n=1}^{i}h_n\cdot\mathbf{k}^4_n\otimes\sigma\cdot a^4=C
\end{equation}

Here \(C\) is the gravitational field (m\(^4\)).

The gravitational source may have spin.

At the same time, the reason why \(\mathbf{k}^4_n\) in this formula does not carry an absolute value sign is that under this system it is default to have a real part and an imaginary part.

The details of \(Z\) inside need to be understood in Chapter 3.

And note that the result of this formula is almost no different from the Einstein field equation, because it also writes out the bending of spacetime under discreteness.

\begin{equation}
h_n=\frac{\chi\rho_n}{b}
\end{equation}

\begin{equation}
b=\Delta_{++}\left(\frac{x^2}{2}\right)=1\ (1/L^2)
\end{equation}

\(\rho_n\) is the energy density at \(n\).

At the same time it is a forward difference on space.

Here \(\rho\) represents the energy density.

Its physical meaning is that a tiny change of mass will cause a tiny change of the discrete variable of spacetime, and the change of the discrete variable will affect the tiny change of spacetime.

The dimension of the left side of this formula and of the right side are the same.

Where on the right side of the equation \(\chi=8\pi G/c^4\).

\begin{equation}
s^2\nabla^2(\mathbf{k})=\frac{\Delta_{2++}(\mathbf{k})}{\Delta_{2++}t^2}
\end{equation}

\(s\) is the propagation speed constant of \(\mathbf{k}\) (dimensionless).

This theory regards \(\mathbf{k}\) as a wave, of course it is also a kind of change of the field.

Because \(c=1\), the second partial derivative with respect to time on the right side of the equation can also become with respect to space.

Its physical meaning is as follows: if \(\mathbf{k}\) has a tiny change, then its propagation with respect to time will propagate in the three directions \(x\), \(y\), \(z\) at the propagation speed.

Because \(\mathbf{k}\) is dimensionless, the dimensions of the left and right sides of this formula are also the same.

And if \(\mathbf{k}\) is complex, then write it in the form of \(e\) again...

\begin{equation}
\mathbf{k}^4a^3(ca_t)
\end{equation}

After computing the rate of change, compute the small volume of one place, and because there is \(a_t\), therefore what is computed is the four-dimensional figure. (You should still remember \(c=1\), space and time are equivalent)


A body in the gravitational field moves along the geodesic of the curved spacetime, i.e. the one that follows the curvature among the two shortest lines in Chapter 3; the part \(s\) that jumps over spacetime is not counted for now.

\section{Z Space (Z is named with the first letter of my name)}

Let us imagine, the \(x\) axis, \(y\) axis, \(z\) axis, \(t\) axis of time and space, all things seem to be spread out on one space, for example, any behavior carried out on a sphere should only be confined to this sphere. Therefore, we call this sphere Z space.

Here \(Z_x\) represents the Z space on the \(x\) axis, that is, countless \(x\) axes make up this figure, of course there is only one \(x\) axis, that is undefined. Although countless lines can also make up a solid, \(x\) still needs to depend on the others, and the other axes are the same by analogy.

Here \(k_x\) is taking \(\Delta_+\) of \(Z_x\).

And to change \(k_x\) back one needs the summation sign:

\begin{equation}
\sum_{n=1}^{x}k_{xn}a=Z_x
\end{equation}

where \(x\) represents how many \(k_{xn}a\) there are in the Z space on \(x\).

You also need to know that at every coordinate, the change at every point is not necessarily the same, spacetime is not always Euclidean geometry.

Imagine a circle and a line, passing by the edge of the circle, and the point where they touch, the change of that \(a\) is \(k_x\), and then many lines are shot over to cover every place of it, note that they are all \(x\) axis.

No boundary is the situation where \(n\), \(n+1\) are undefined. Take the simplest sphere as an example, it is a four-dimensional sphere, then it has no boundary, because even if we walk to its end, one more step forward is only returning to the starting point.

And those unfolding pictures of the universe are essentially transforming 3+1 dimensions into 2+1 dimensions, therefore, after unfolding it is not the original one, because it went from having no boundary to having a boundary.

If there is purely only change on the \(z\) axis, it shows that space only has the \(z\) axis, that is, one dimension.

And the complete Z space should be all the axes contained within this dimension.

At the same time here, I have a conjecture: between two points the shortest line segment has two, except in absolute vacuum.

Because it is not absolute vacuum, therefore there will certainly be energy, therefore spacetime will also bend; suppose before bending there are point \(a\) and point \(b\), they are influenced by the bending, that is to say, from point \(a\) to point \(b\) there is the straight line that jumps over the bent spacetime, and there is also the one that follows the bent spacetime, and these two are indeed both shortest.

But these two segments have one common point, that is, you cannot find the negative segment corresponding to them to cancel; we call the part of these two segments that jumps over spacetime \(s\), then its energy cannot be simply the known segment multiplied by some constant, but should be the sum of the known segment and \(s\) multiplied by some constant.

\section{Matrix with \(i\)}

Here, we assume one situation of one point: and it can only be one point

\begin{align*}
k_x &= 3+k_yi\\
k_y &= 24+k_x\\
k_z &= e^i\\
k_t &= e^i
\end{align*}

\(\mathbf{k}^4\) written in the form of a matrix:

\begin{equation}
\mathbf{k}^4=
\begin{pmatrix}
3 & i & 0 & 0\\
1 & 24 & 0 & 0\\
0 & 0 & e^i & 0\\
0 & 0 & 0 & e^i
\end{pmatrix}
\end{equation}

This matrix is non-Hermitian (remember to add).

Here, I need to state one point first, why I use forward difference and not backward difference: first, backward difference fails at the starting point, and forward difference fails at the boundary; if according to the starting point \(x-a\) equals the boundary \(n\), and the starting point \(n+a\) equals the starting point \(x\), then the difference between these two is actually very small.

\section{Weak force}

Currently cannot write it out.

\section{Electromagnetic force}

Same as the weak force.

\section{The minimal action in 3+1 dimensions}

\begin{equation}
\mathcal{Z}=\operatorname{Re}\left(\sum_i\left[\frac{1}{2}\left(\Delta_t\mathbf{k}\right)^2p^2-\frac{1}{2}\mathbf{k}^2p^2s^2\right]a^4
-\sum_i\sum_n V(a_r)\cdot a^4\cdot\mathbf{k}^4_n
+\sum_i\sum_n\frac{\chi\rho_n}{b}\cdot a^4\cdot\mathbf{k}^4_n\right)
\end{equation}

\(p\) is a constant, \(s\) has been mentioned before. At the same time, the \(\mathbf{k}\) in the minuend and in the subtrahend are both dimensionless; \(\Delta_t\mathbf{k}\) represents the change of \(\mathbf{k}\) in one time interval, it is not a forward difference.

\section{Nuclear force}

Before this we need to understand several basic conjectures of ZG:

\begin{enumerate}
\item The distance between any object and another object is at minimum \(a\).
\item If the distance between two objects is smaller than \(a\), or a force wants to make the distance smaller than \(a\), then there will be an opposite force to resist.
\end{enumerate}

\(a\) is only that we take the Planck distance, but it is not fixed.

From this we obtain the formula:

\begin{equation}
V(a_r)=V_{\omega}+V_{2\pi}+V_{ten}+V_{ls}
\end{equation}

\begin{equation}
V_{\omega}=w_{\omega}\frac{e^{-a_r/R_{\omega}}}{a_r},\qquad
V_{2\pi}=-w_{2\pi}\frac{e^{-a_r/R_{2\pi}}}{a_r}
\end{equation}

\begin{equation}
V_{ten}=w_T\,S_{12}\frac{1}{x}\left(1+\frac{3}{x}+\frac{3}{x^2}\right)e^{-x}\left[1-e^{-(a_r/R_{\omega})^2}\right],\qquad x=\frac{a_r}{R_{\pi}}
\end{equation}

\begin{equation}
V_{ls}=w_V\,(\mathbf{L}\cdot\mathbf{S})\,e^{-a_r/R_{2\pi}}
\end{equation}

\begin{equation}
R_{\omega}=\frac{\hbar}{m_{\omega}c}=0.25\ \mathrm{fm},\quad
R_{2\pi}=\frac{\hbar}{2m_{\pi}c}=0.70\ \mathrm{fm},\quad
R_{\pi}=\frac{\hbar}{m_{\pi}c}=1.41\ \mathrm{fm}
\end{equation}

\begin{equation}
S_{12}=3(\sigma_1\cdot\hat a_{rr})(\sigma_2\cdot\hat a_{rr})-(\sigma_1\cdot\sigma_2),\qquad \hat a_{rr}=\frac{a_r\ (\text{vector})}{a_r\ (\text{value})}
\end{equation}

Here \(r\) is taken as \(a\), so that this can hold in comparison with present real results.

At the same time one needs to know, \(\mathbf{k}^4\) has changed, so \(a\) needs to multiply \(\mathbf{k}^4\), and will also change.

Here the strengths are constants: \(w_{\omega}\approx 2842.92\ \mathrm{MeV\cdot fm}\), \(w_{2\pi}\approx 479.83\ \mathrm{MeV\cdot fm}\) (divided by length, hence the fm), \(w_T\approx -9.195\ \mathrm{MeV}\), \(w_V\approx 39.676\ \mathrm{MeV}\) (multiplied by dimensionless shapes, hence no fm).

\(a_r\) is the distance between two objects where nuclear force happens.

Using this can explain nuclear force.


\subsection{Motion of a nucleon in the field}

\begin{equation}
\Delta_{++}x=-\frac{\Delta_+V(a_r)}{m}
\end{equation}


A numerical comparison of four deuteron observables is carried out by ZGFT\_simplex.py
(Nelder--Mead simplex, residual $\chi^2=0.0004$):

\begin{center}
\begin{tabular}{lccc}
\hline
observable & theory & experiment & deviation \\
\hline
deuteron energy $E_0$ (MeV) & $-2.2381$ & $-2.224$ & $-0.63\%$ \\
triplet scattering length $a_t$ (fm) & $5.445$ & $5.42$ & $+0.46\%$ \\
D-wave probability $P_D$ & $5.694\%$ & $5.8\%$ & $-1.83\%$ \\
quadrupole moment $Q$ (fm$^2$) & $0.28652$ & $0.2859$ & $+0.22\%$ \\
\hline
\end{tabular}
\end{center}

\section{Lagrangian density at this stage}

\begin{equation}
L=\operatorname{Re}\left(\frac{1}{2}\Delta_t\mathbf{k}^2p^2-\frac{1}{2}\mathbf{k}^2p^2s^2-V(a_r)+\frac{\chi\rho}{b}\right)
\end{equation}

\begin{equation}
\mathcal{Z}=\operatorname{Re}\left(\sum_i\sum_n L\cdot a^4\cdot\mathbf{k}^4_n\right)
\end{equation}



\section{To be completed}

\textbf{$\gamma$ matrices and the exact coefficients of $w_T$, $w_V$}

$w_T$ and $w_V$ are presently independent parameters. If they come from $\omega$ vector exchange, they should be

\begin{equation}
w_T=w_{\omega}\frac{(\hbar/2Mc)^2}{R_{\omega}^2}c_T,\qquad
w_V=w_{\omega}\frac{(\hbar/2Mc)^2}{R_{\omega}^2}c_l
\end{equation}

where $c_T$, $c_l$ are fixed by $\gamma$-matrix elements, requiring the Dirac equation and $\gamma$-matrix algebra.


\vspace{2em}
\noindent\textbf{Declaration}

The core ideas and core formulas of this work were proposed by me (Ziqi Zhi); AI only assisted in
refining the derivations, calculations and typesetting. This work is made by me, with thanks to DeepSeek AI.

\vspace{1em}
\noindent Ziqi Zhi\quad 2026

\end{document}
